%% 2i2d-cycle-vie-materiau — les cinq étapes du cycle de vie d'un matériau, avec
%% l'impact environnemental de chacune ; en fin de vie, trois voies de valorisation,
%% dont deux referment la boucle (réutilisation → fabrication, recyclage → extraction).
%% La station « Fin de vie » est placée à droite du cercle pour que les trois voies
%% partent vers l'extérieur sans traverser le cycle.
\documentclass[tikz,border=4pt]{standalone}
\usepackage{liaisons}
\begin{document}
\begin{tikzpicture}[x=1cm,y=1cm,
  station/.style={draw=solideB, fill=solideB!8, line width=0.9pt, rounded corners=3pt,
                  text width=1.68cm, align=center, inner sep=2.5pt},
  boucle/.style={-{Stealth[length=5pt]}, line width=1pt, solideD},
  voie/.style={-{Stealth[length=5pt]}, line width=0.9pt, solideE},
  retour/.style={-{Stealth[length=5pt]}, line width=0.8pt, solideE,
                 dash pattern=on 3pt off 2pt},
  sortie/.style={draw=solideE, fill=solideE!8, line width=0.8pt, rounded corners=2.5pt,
                 inner sep=3pt, align=left},
  petit/.style={font=\scriptsize, inner sep=1.5pt},
  retlab/.style={font=\scriptsize, text=solideE, fill=white, inner sep=1.5pt}]

%% --- gabarit d'une station : nom en gras, impact en petit gris ------------------
\newcommand\etape[2]{{\scriptsize\bfseries\color{solideB} #1}\\[1pt]{\tiny\color{black!55} #2}}

%% --- les cinq stations, sur un cercle parcouru dans le sens horaire -------------
\def\cx{2.3} \def\cy{2.5} \def\rr{1.75}
\newcommand\pos[1]{({\cx+\rr*cos(288-72*#1)},{\cy+\rr*sin(288-72*#1)})}
\node[station] (S1) at \pos{0} {\etape{Extraction}{carrières, forages}};
\node[station] (S2) at \pos{1} {\etape{Transport}{GES}};
\node[station] (S3) at \pos{2} {\etape{Fabrication}{GES}};
\node[station] (S4) at \pos{3} {\etape{Utilisation}{GES}};
\node[station] (S5) at \pos{4} {\etape{Fin de vie}{déchets, pollution}};

%% --- le cycle, sens horaire ------------------------------------------------------
\foreach \a/\b in {S1/S2, S2/S3, S3/S4, S4/S5, S5/S1} { \draw[boucle] (\a) to[bend right=14] (\b); }
\node[petit, text=solideD, align=center] at (\cx,\cy) {sens\\du cycle};

%% --- les trois voies de valorisation, à droite ------------------------------------
\node[sortie, anchor=west] (V1) at (5.20,4.05)
     {{\scriptsize\bfseries Réutilisation}\\[1pt]{\scriptsize l'objet ressert}};
\node[sortie, anchor=west] (V2) at (5.20,2.50)
     {{\scriptsize\bfseries Recyclage}\\[1pt]{\scriptsize la matière ressert}};
\node[sortie, anchor=west] (V3) at (5.20,0.95)
     {{\scriptsize\bfseries Valorisation énergétique}\\[1pt]{\scriptsize la chaleur est récupérée}};
\draw[voie] (S5.north east) to[out=55,in=180]  (V1.west);
\draw[voie] (S5.east)       to[out=0,in=180]   (V2.west);
\draw[voie] (S5.south east) to[out=-55,in=180] (V3.west);

%% --- les deux retours qui referment la boucle ---------------------------------------
\draw[retour] (V1.north) to[out=90,in=75,looseness=1.25]
      node[retlab, pos=0.55] {retour vers la fabrication} (S3.north east);
\draw[retour] (V2.south west) to[out=-125,in=20,looseness=0.9]
      node[retlab, pos=0.4, anchor=south] {retour vers l'extraction} (S1.east);

%% --- note de lecture ------------------------------------------------------------------
\node[petit, anchor=north west, text=black!55, font=\scriptsize\itshape] at (0.0,0.20)
     {GES : gaz à effet de serre};
\end{tikzpicture}
\end{document}
